\chapter{Design Requirements and Evaluation Scope}

Nebula represents a target using minimum locked-phase eye height, minimum eye
width, minimum decision margin, and maximum CTLE power. Peaking, HD3, integrated
noise, and observed DFE errors are additional evaluator gates when their stages
are enabled. These are project-defined engineering targets unless an external
requirement is explicitly cited.

\begin{table}[htbp]
\centering
\caption{Principal receiver targets and their precise implementation scope.}
\label{tab:requirements}
\begin{tabularx}{\textwidth}{@{}l S[table-format=2.3] l X@{}}
\toprule
Metric & {Target} & Unit & Interpretation \\
\midrule
Locked-phase DFE eye height & 0.1 & V & Lower bound on the behavioral DFE metric. \\
DFE eye width & 0.4 & UI & Lower bound over the simulated finite bit pattern. \\
DFE decision margin & 0.0 & V & Corrected sample margin at the locked phase. \\
CTLE power & 0.015 & W & Upper bound for CTLE supply power, not total receiver power. \\
HD3 & -30 & dB & Upper bound for the implemented \SI{100}{MHz} tone test. \\
Input-referred noise & 1.5 & mV RMS & Upper bound integrated from \SI{10}{MHz} to \SI{5}{GHz}. \\
\bottomrule
\end{tabularx}
\end{table}

The phrases ``eye height'' and ``power'' are otherwise ambiguous. This report
therefore uses exact names such as \emph{locked-phase DFE eye height} and
\emph{CTLE power}. Zero observed errors always refers to the evaluated finite
sequence; it is not a bit-error-rate guarantee.

\section{Qualification sets}

The implementation defines a 60-condition grid from five process corners
(TT, SS, FF, SF, and FS), three supply voltages, and four temperatures. The
primary completed run instead evaluates 27 conditions: TT/SS/FF,
\SIlist{1.71;1.80;1.89}{V}, and \SIlist{0;27;125}{\degreeCelsius}. It omits
\SI{75}{\degreeCelsius}, SF, and FS. Throughout this report, ``demonstrated
reduced 27-condition grid'' names this recorded result; ``extended 60-condition
grid'' names the uncompleted competition-aligned scope.
